
\section{Introduction}

\subsection{Motivation}

Transformer language models concentrate the majority of their parameters in 
the feed-forward network and the multi-head self-attention (MHSA) mechanism. 
In DistilBERT-base, the four attention projections $W_Q, W_K, W_V, W_O$ 
account for roughly \textbf{70\%} of the parameters in each transformer block. 
Reducing the parameter cost of attention is a high-leverage compression 
strategy.

Standard MHSA stores four dense matrices per layer:

\[
W_Q, W_K, W_V, W_O \in \mathbb{R}^{d \times d}
\tag{1.1}
\]

totaling $4 d^2$ parameters. For DistilBERT ($d = 768$), this is 
$4 \times 589{,}824 = 2{,}359{,}296$ parameters per layer, or 14.2M across 
six layers.

\textbf{Low-rank factorization} is the classical primitive for reducing this cost. 
For each $W$, one replaces it with a product:

\[
W \approx U V, \qquad
U \in \mathbb{R}^{d_\text{out} \times r}, \quad
V \in \mathbb{R}^{r \times d_\text{in}}, \quad r \ll \min(d_\text{in}, d_\text{out})
\tag{1.2}
\]

The parameter count drops from $d_\text{in} d_\text{out}$ to 
$r(d_\text{in} + d_\text{out})$. For square matrices, the reduction factor is 
$d / 2r$. Classical SVD (Eckart--Young, 1936) gives the optimal rank-$r$ 
approximation in Frobenius norm.

\textbf{The rank-selection problem.} Classical SVD requires choosing $r$ 
a-priori, which is a hard trade-off:

\begin{itemize}
\item \textbf{Too small} -- approximation error is high, downstream accuracy drops.
\item \textbf{Too large} -- parameter reduction is minimal, defeating the purpose.

\end{itemize}
In practice, different layers have different \textbf{intrinsic ranks}. Attention 
projections in early layers of DistilBERT tend to have rapidly-decaying 
spectra (effective rank ~30--48), while middle layers decay more slowly 
(effective rank ~80--100). A single global $r$ is suboptimal for both.

\subsection{Problem Statement}

Given a pretrained transformer $\mathcal{M}$ with $N_l$ layers and a target 
energy threshold $\eta \in (0, 1)$, we seek a compressed model 
$\mathcal{M}'$ such that:

\[
\mathcal{M}' \approx \mathcal{M}, \qquad
\frac{|\theta(\mathcal{M}')|}{|\theta(\mathcal{M})|} \ll 1
\tag{1.3}
\]

subject to:

\begin{enumerate}
\item \textbf{Per-layer adaptive rank.} Different layers may use different
  $r_\ell$.
\item \textbf{Spectral energy preservation.} $E(r_\ell) \geq \eta$ for each layer.
\item \textbf{Parameter reduction.} $r_\ell(d_\text{in} + d_\text{out}) < d_\text{in} d_\text{out}$.
\item \textbf{Multi-architecture support.} Flat, nested, causal, and GQA transformers.
\item \textbf{No base retraining.} Only LoRA adapters may be trained.

\end{enumerate}
\subsection{Our Approach: MinimalSVD}

\textbf{MinimalSVD} is a compact, auditable baseline for low-rank attention 
compression. It has two components:

\begin{enumerate}
\item \textbf{SVD factorization with energy-adaptive rank.} Each $W$ is factorized as
  $U V$ via truncated SVD. The rank is selected per layer as:

  $$
  r\_\ell = \min\!\left(r\_\text{max}, \left|\{i : E(i) < \eta\}\right| + 1\right)
  \tag{1.4}
  $$

  where $E(r) = \sum_{i \leq r} \sigma_i^2 / \sum_i \sigma_i^2$. The
  factorization distributes $\Sigma^{1/2}$ symmetrically.

\item \textbf{LoRA recovery.} After factorization, base weights are frozen and LoRA
  adapters $A, B$ are added with $B = 0$ at initialization.

\end{enumerate}
\textbf{Design philosophy.} MinimalSVD intentionally keeps the injection protocol 
\textbf{minimal} -- the injection hook targets a single transformer layer at a 
time, delegating whole-model traversal to the caller. This makes the code 
easy to audit, extend, and embed in different frameworks. Unlike the more 
complex variants (NOVA, PRISM), MinimalSVD does not perform pre-testing on 
the first layer; correctness is ensured by the strict weight-extraction 
checks in \texttt{copy\_weights}.

\subsection{Contributions}

\begin{itemize}
\item \textbf{A minimal reference implementation.} A compact, auditable baseline for
  energy-adaptive low-rank attention compression.
\item \textbf{Spectral energy-adaptive rank selection.} A principled criterion that
  selects a different rank per layer based on spectral decay.
\item \textbf{Multi-architecture support.} Four families (DistilBERT, RoBERTa, GPT-2,
  Qwen3.5) with architecture-aware weight extraction.
\item \textbf{A systematic evaluation.} Benchmark on three models on SST-2 with
  ablations over $r_\text{max}$, $\eta$, and LoRA rank.
\item \textbf{Full reproducibility.} Code, checkpoints, and scripts released under
  Apache 2.0.

\end{itemize}
\subsection{Relationship to Companion Variants}

MinimalSVD is the \textbf{baseline} of a family of low-rank attention modules 
that share the same \texttt{FactorisedLinear} core:

\begin{table}[H]
\centering
\small
\begin{tabular}{ll}
\toprule
\textbf{Variant} & \textbf{Additional features} \\
\midrule
\textbf{MinimalSVD} & Baseline (this work) \\
MinimalNOVA & Whole-model injection + pre-test \\
MinimalPRISM & Injection-time nested detection \\
\bottomrule
\end{tabular}
\end{table}

All three share the same energy-adaptive rank selection, LoRA recovery, and 
weight transfer logic. MinimalSVD exists to isolate the core contribution 
from engineering additions.

\subsection{Paper Organization}

Section 2 reviews related work. Section 3 provides background. Section 4 
presents MinimalSVD. Section 5 describes architecture integration. Section 
6 details the setup. Section 7 reports results. Section 8 analyzes the 
energy-adaptive rank behavior. Section 9 discusses guidance. Section 10 
lists limitations. Section 11 concludes.



\subsection{Notation}

\begin{lstlisting}

---
\end{lstlisting}